% Replace the placeholder fields below before final submission.
\newcommand{\dissertationtitle}{One Size Does Not Fit All: Personal Match and Marketing Message Effectiveness on Social Networks}
\newcommand{\studentname}{WANG, Yanjun}
\newcommand{\studentnumber}{3036492833}
\newcommand{\supervisorname}{Prof. Max Z.J. Shen}
\newcommand{\submissiondate}{2\textsuperscript{nd} August 2026}

\pagenumbering{gobble}

\begin{titlepage}
\thispagestyle{empty}
\begin{singlespace}
\begin{center}
{\bfseries\fontsize{20}{24}\selectfont The University of Hong Kong\par}
\vspace{0.35cm}
{\fontsize{16}{19.2}\selectfont Department of Data and Systems Engineering\par}

\vspace{1.05cm}

{\fontsize{16}{19.2}\selectfont M.Sc.(Eng.) in Industrial Engineering and Logistics\\
Management\par}

\vspace{1.35cm}

{\itshape\fontsize{16}{19.2}\selectfont DASE 7099 Dissertation\par}

\vspace{1.1cm}

{\itshape\fontsize{16}{19.2}\selectfont Final Report\par}
\vspace{-0.25cm}
{\rule{3.2cm}{0.4pt}\par}

\vspace{1.45cm}

\begin{minipage}{0.82\textwidth}
\centering
{\fontsize{14}{16.8}\selectfont \dissertationtitle\par}
\end{minipage}

\vspace{1.45cm}

{\fontsize{14}{16.8}\selectfont \studentname\ (\studentnumber)\par}

\vspace{1.25cm}

{\bfseries\fontsize{14}{16.8}\selectfont Supervisor:\par}
\vspace{0.2cm}
{\fontsize{14}{16.8}\selectfont \supervisorname\par}

\vfill

{\fontsize{14}{16.8}\selectfont \textbf{Date of submission:} \submissiondate\par}
\end{center}
\end{singlespace}
\end{titlepage}

\clearpage
\thispagestyle{empty}
\frontmatterheading{Abstract}

In social media marketing, firms increasingly rely on employees and marketing agents to disseminate promotional content through their own interpersonal networks. Yet marketing messages do not diffuse equally well across all agents, which points to a limitation of one-size-fits-all content distribution. This dissertation examines how article content, observable agent characteristics, and content–agent fit are associated with the diffusion effectiveness of marketing messages, using real-world WeChat marketing log data generated by a single interior design company on the WeiMai social marketing platform. Diffusion effectiveness is treated as a multi-dimensional concept and evaluated in reach, cascade depth, resharing, and structural virality, among others.

The empirical design combines LLM(Large Language Model)–assisted content feature extraction, diffusion cascade reconstruction, and a leakage-aware three-level analytical framework. Content features are constructed from article titles, body text, and image-use metadata, and include topic labels, topic probability scores, image indicators, word count, and a title–content semantic consistency measure. Diffusion cascades are reconstructed from sharing linkage records to correct layer assignments and derive multi-dimensional diffusion outcomes. The analysis separates three levels with distinct units of analysis: content-level article–agent cases (6,557 cases), agent-characteristic-level records (592 agents), and content–agent matching-level agent–topic pairs (1,142 pairs). Keeping the levels separate avoids combining different units of analysis in a single model and preserves a clear interpretation of each coefficient.

The results indicate that, at the content level, title–content semantic consistency is positively associated with diffusion reach. At the agent level, observable agent characteristics relate to some diffusion dimensions. Job category corresponds to systematic differences across several dimensions. Marketing and operations agents, for instance, show larger and deeper cascades than field-sales agents but occupy lower network-centrality positions. At the agent–topic level, a stronger match between an article's topics and the sharing agent's professional context is positively associated with diffusion performance, and this association is also robust to a binary specification of the match measure.

These findings suggest that social marketing strategies should move beyond uniform content distribution towards personalised content assignment, with particular attention to the fit between article topics and agents' professional contexts. The study contributes a framework that integrates LLM-assisted content analysis, corrected diffusion-cascade reconstruction, and three-level empirical modelling and analysis of multiple diffusion outcomes. The analysis is explanatory and associational. The findings indicate systematic relationships but do not identify causal effects.

\clearpage
\pagenumbering{roman}
\setcounter{page}{1}

\phantomsection
\addcontentsline{toc}{chapter}{Declaration}
\frontmatterheading{Declaration}

\begin{doublespace}
I, \studentname, hereby declare that the M.Sc.(Eng.) Dissertation final report, entitled ``\dissertationtitle'', which I am submitting, represents my own work and has not been previously submitted to this or any other institutions in the application for admission to a degree, a diploma or any other qualifications.
\end{doublespace}

\vspace{2.4cm}

\begin{flushright}
\begin{tabular}{@{}r@{\hspace{0.8em}}p{6.5cm}@{}}
Signed: & \underline{\makebox[6.5cm][l]{\studentname}} \\[1.0em]
Date:   & \underline{\makebox[6.5cm][l]{\submissiondate}}
\end{tabular}
\end{flushright}

\clearpage
\phantomsection
\addcontentsline{toc}{chapter}{Table of Contents}
\fronttableofcontents

\clearpage
\phantomsection
\addcontentsline{toc}{chapter}{List of Figures}
\frontlistoffigures

\clearpage
\phantomsection
\addcontentsline{toc}{chapter}{List of Tables}
\frontlistoftables

\clearpage
